\chapter{Непрекъснати съвместни разпределения, Гама разпределение}

\section{Съвместни непрекъснати разпределения}

\begin{defbox}[~(Съвместно непрекъснато разпределение)]
Нека $X=(X_1,\dots,X_n)$ е вектор от $n$ случайни величини.
Казваме, че $X$ има \textbf{съвместно непрекъснато разпределение} със
\textbf{съвместна плътност} $f_X:\mathbb R^n\to\mathbb R$, ако:
\begin{enumerate}[label=\arabic*)]
  \item $f_X(x_1,\dots,x_n)\ge0$;
  \item $\displaystyle\int_{\mathbb R^n} f_X(x_1,\dots,x_n)\,dx_1\cdots dx_n = 1$;
  \item за всяко $D\subseteq\mathbb R^n$: $\displaystyle\mathbb P(X\in D) = \int_D f_X(x_1,\dots,x_n)\,dx_1\cdots dx_n$.
\end{enumerate}
\end{defbox}

\begin{defbox}[~(Носител)]
$D_{f_X} = \{\bar x\in\mathbb R^n : f_X(\bar x)>0\}$.
\end{defbox}

\begin{defbox}[~(Маргинална плътност)]
Нека $X=(X_1,\dots,X_n)$ има съвместна плътност $f_X$. Тогава
\[
  f_{X_j}(x_j) = \int_{-\infty}^{\infty}\cdots\int_{-\infty}^{\infty} f_X(x_1,\dots,x_{j-1},x_j,x_{j+1},\dots,x_n)\,dx_1\cdots dx_{j-1}\,dx_{j+1}\cdots dx_n
\]
(интегрираме по всички променливи с~изключение на $x_j$, който
остава фиксиран) се нарича \textbf{маргинална плътност} на $X_j$.
\end{defbox}

\begin{defbox}[~(Условна плътност)]
Нека $X=(X_1,X_2)$ със съвместна плътност $f_X$ и~нека $x_2$ е
фиксирано с~$f_{X_2}(x_2)>0$. Тогава
\[
  f_{X_1\mid X_2=x_2}(x_1\mid x_2) := \frac{f_X(x_1,x_2)}{f_{X_2}(x_2)}
\]
е \textbf{условната плътност} на $X_1$ при информация за $X_2=x_2$.
За нея $\displaystyle\mathbb P\big(X_1\in(a,b)\mid X_2=x_2\big) =
\int_a^b f_{X_1\mid X_2=x_2}(x_1\mid x_2)\,dx_1$.
\end{defbox}

\begin{defbox}[~(Съвместна функция на разпределение)]
Нека $X=(X_1,\dots,X_n)$. Тогава
\[
  F_X(x_1,\dots,x_n) = \mathbb P(X_1<x_1,\dots,X_n<x_n), \qquad \forall(x_1,\dots,x_n)\in\mathbb R^n.
\]
Ако $X$ е вектор от непрекъснати случайни величини със съвместна
плътност $f_X$, непрекъсната в~$\bar x\in\mathbb R^n$, то
\[
  f_X(\bar x) = \frac{\partial^n F_X(\bar x)}{\partial x_1\partial x_2\cdots\partial x_n}, \qquad
  F_X(\bar x) = \int_{-\infty}^{x_1}\!\!\int_{-\infty}^{x_2}\!\!\cdots\!\!\int_{-\infty}^{x_n} f_X(y_1,\dots,y_n)\,dy_1\cdots dy_n.
\]
\end{defbox}

\begin{defbox}[~(Независимост)]
Нека $X=(X_1,X_2)$. Казваме, че $X_1\perp\!\!\!\perp X_2$, ако
\[
  \forall(x_1,x_2)\in\mathbb R^2:\quad F_X(x_1,x_2) = F_{X_1}(x_1)F_{X_2}(x_2)
\]
\[
  \big(\text{т.е.\ } \mathbb P(X_1<x_1,X_2<x_2)=\mathbb P(X_1<x_1)\mathbb P(X_2<x_2)\big),
\]
еквивалентно на $f_X(x_1,x_2) = f_{X_1}(x_1)\cdot f_{X_2}(x_2)$ -
съвместната плътност се разпада на произведение от двете маргинални
плътности.
\end{defbox}

\begin{defbox}[~(Независимост в~съвкупност)]
Нека $X=(X_1,\dots,X_n)$ има съвместна плътност $f_X$. Казваме, че
$X_1,\dots,X_n$ са \textbf{независими в~съвкупност}, ако за всяко
$\bar x=(x_1,\dots,x_n)\in\mathbb R^n$:
\[
  f_X(x_1,\dots,x_n) = \prod_{j=1}^{n} f_{X_j}(x_j).
\]
\end{defbox}

\section{Очакване за случаен вектор}

\begin{rembox}[Смяна на променливите за очакване (много измерения)]
Нека $X$ е вектор от случайни величини със съвместна плътност $f_X$ и~$g:\mathbb R^n\to\mathbb R$. Тогава
$\displaystyle \mathbb Eg(X) = \int_{\bar x\in\mathbb R^n} g(\bar x)\,f_X(\bar x)\,d\bar x$.
\end{rembox}

\begin{thmbox}[]{Теорема}
Нека $X=(X_1,X_2)$ е непрекъснат случаен вектор със съвместна
плътност $f_X$ и~нека $\mathbb EX_1,\mathbb EX_2$ съществуват.
Тогава $\mathbb E(X_1+X_2) = \mathbb EX_1+\mathbb EX_2$. \emph{(Тази
формула важи за всякакви случайни величини - дискретни, непрекъснати,
независими или не.)}
\end{thmbox}

\begin{proofbox}
Нека $g(x_1,x_2)=x_1+x_2$. Тогава
\begin{align*}
  \mathbb E(X_1+X_2) &= \iint (x_1+x_2)f_X(x_1,x_2)\,dx_1dx_2\\
  &= \int x_1\Big(\int f_X(x_1,x_2)\,dx_2\Big)dx_1 + \int x_2\Big(\int f_X(x_1,x_2)\,dx_1\Big)dx_2
\end{align*}
\[
  = \int x_1 f_{X_1}(x_1)\,dx_1 + \int x_2 f_{X_2}(x_2)\,dx_2 = \mathbb EX_1+\mathbb EX_2. \qedbg
\]
\end{proofbox}

\begin{thmbox}[]{Теорема}
Нека $X=(X_1,X_2)$ е непрекъснат случаен вектор, $X_1\perp\!\!\!\perp
X_2$, и~нека $DX_1,DX_2$ съществуват. Тогава $D(X_1+X_2) = DX_1+DX_2$.
\end{thmbox}

\begin{proofbox}
\begin{align*}
  D(X_1+X_2) &= \mathbb E\big[(X_1+X_2)-\mathbb E(X_1+X_2)\big]^2 = \mathbb E\big[(X_1-\mathbb EX_1)+(X_2-\mathbb EX_2)\big]^2\\
  &= \mathbb E(X_1-\mathbb EX_1)^2 + \mathbb E(X_2-\mathbb EX_2)^2 + 2\,\mathbb E(X_1-\mathbb EX_1)(X_2-\mathbb EX_2)\\
  &= DX_1+DX_2+2\,\mathbb E(X_1-\mathbb EX_1)(X_2-\mathbb EX_2).
\end{align*}
Остава да покажем, че последният член е нула. Понеже $X_1\perp\!\!\!
\perp X_2$, съвместната плътност се разпада: $f_X(x_1,x_2)=
f_{X_1}(x_1)f_{X_2}(x_2)$. Тогава
\[
  \mathbb E(X_1-\mathbb EX_1)(X_2-\mathbb EX_2) = \iint (x_1-\mathbb EX_1)(x_2-\mathbb EX_2)f_{X_1}(x_1)f_{X_2}(x_2)\,dx_1dx_2
\]
\[
  = \Big(\int (x_1-\mathbb EX_1)f_{X_1}(x_1)\,dx_1\Big)\Big(\int (x_2-\mathbb EX_2)f_{X_2}(x_2)\,dx_2\Big),
\]
а~всеки от двата множителя е нула: $\int(x_2-\mathbb EX_2)f_{X_2}
(x_2)\,dx_2 = \int x_2 f_{X_2}(x_2)dx_2 - \mathbb EX_2\int f_{X_2}
(x_2)dx_2 = \mathbb EX_2 - \mathbb EX_2\cdot1 = 0$ (аналогично за
първия множител). Следователно $D(X_1+X_2)=DX_1+DX_2$. \qedbg
\end{proofbox}

\begin{thmbox}[]{Теорема}
Нека $X=(X_1,X_2)$ е непрекъснат случаен вектор, $X_1\perp\!\!\!\perp
X_2$, и~$\mathbb EX_1,\mathbb EX_2$ съществуват. Тогава $\mathbb
EX_1X_2 = \mathbb EX_1\cdot\mathbb EX_2$.
\end{thmbox}

\begin{proofbox}
\[
  \mathbb EX_1X_2 = \iint x_1x_2\,f_X(x_1,x_2)\,dx_1dx_2 \overset{X_1\perp\!\!\!\perp X_2}{=} \iint x_1x_2\,f_{X_1}(x_1)f_{X_2}(x_2)\,dx_1dx_2
\]
\[
  = \Big(\int x_1 f_{X_1}(x_1)\,dx_1\Big)\Big(\int x_2 f_{X_2}(x_2)\,dx_2\Big) = \mathbb EX_1\cdot\mathbb EX_2. \qedbg
\]
\end{proofbox}

\section{Смяна на променливите (многомерна)}

\begin{thmbox}[]{Теорема}
Нека $X=(X_1,X_2)$ е случаен вектор със съвместна плътност $f_X$ и~носител $D_{f_X}\subseteq\mathbb R^2$, и~нека $g:D_{f_X}\to\mathbb
R^2$, $Y=g(X)=(g_1(X_1,X_2),g_2(X_1,X_2))$, е биективна. Нека
$h=g^{-1}$, $X=h(Y)=(h_1(Y_1,Y_2),h_2(Y_1,Y_2))$, като $g$ и~$h$ са
непрекъснати и~диференцируеми в~$D_{f_X}$ и~$g(D_{f_X})$
съответно. Ако Якобианът
\[
  J(y) = \begin{pmatrix}\dfrac{\partial h_1}{\partial y_1} & \dfrac{\partial h_1}{\partial y_2}\\[1ex] \dfrac{\partial h_2}{\partial y_1} & \dfrac{\partial h_2}{\partial y_2}\end{pmatrix}
\]
удовлетворява $\det J(y)\neq0$ за всяко $y\in g(D_{f_X})$, то $Y$ е
непрекъснат случаен вектор с~\[
  f_Y(y_1,y_2) = f_X\big(h_1(y_1,y_2),h_2(y_1,y_2)\big)\cdot\big|\det J(y)\big|.
\]
\end{thmbox}

\begin{proofbox}
Нека $A\subseteq g(D_{f_X})$. Тогава
\[
  \mathbb P(Y\in A) = \mathbb P\big(g(X)\in A\big) = \mathbb P\big(X\in h(A)\big) = \int_{x\in h(A)} f_X(x)\,dx.
\]
Извършваме смяната на променливите $x=h(y)$ (с~якобиан $J(y)$):
\[
  \int_{x\in h(A)} f_X(x)\,dx = \int_{y\in A} f_X\big(h(y)\big)\cdot|\det J(y)|\,dy,
\]
откъдето, по дефиниция на плътност, $f_Y(y) = f_X(h(y))\cdot|\det J(y)|$. \qedbg
\end{proofbox}

\begin{thmbox}[~(сума на независими нормални случайни величини)]{Теорема}
Нека $X_i\sim N(\mu_i,\sigma_i^2)$, $1\le i\le n$, са независими в~съвкупност. Тогава
\[
  Y = \sum_{i=1}^{n} X_i \sim N\Big(\sum_{i=1}^n \mu_i,\ \sum_{i=1}^n \sigma_i^2\Big).
\]
\end{thmbox}

\begin{proofbox}
Достатъчно е да докажем твърдението за $n=2$ - общият случай следва
по индукция (сумираме последователно). Записваме $X_1=\mu_1+
\sigma_1 V_1$, $X_2=\mu_2+\sigma_2V_2$, където $V_1,V_2\sim N(0,1)$
независими (Тема~10). Тогава
\[
  X_1+X_2 = (\mu_1+\mu_2) + \underbrace{\sigma_1V_1}_{=:W_1\sim N(0,\sigma_1^2)} + \underbrace{\sigma_2V_2}_{=:W_2\sim N(0,\sigma_2^2)},
\]
и~е достатъчно да докажем, че $W_1+W_2\sim N(0,\sigma_1^2+\sigma_2^2)$
за независими $W_1\sim N(0,\sigma_1^2)$, $W_2\sim N(0,\sigma_2^2)$
(тогава $Y=(\mu_1+\mu_2)+(W_1+W_2)\sim N(\mu_1+\mu_2,\sigma_1^2+
\sigma_2^2)$ по теоремата за линейно преобразувание на нормално
разпределение).

Правим трансформацията $W=X_2$, $Z=X_1+X_2$ (т.е.\ $X_1=Z-W$,
$X_2=W$), с~Якобиан
$|J|(z,w)=\left|\det\begin{pmatrix}\partial X_1/\partial Z & \partial X_1/\partial W\\ \partial X_2/\partial Z & \partial X_2/\partial W\end{pmatrix}\right|=\left|\det\begin{pmatrix}1&-1\\0&1\end{pmatrix}\right|=1$.
Понеже $X_1\perp\!\!\!\perp X_2$, съвместната им плътност е
$f_{X_1,X_2}(x_1,x_2)=\dfrac1{2\pi\sigma_1\sigma_2}e^{-\frac{x_1^2}{2\sigma_1^2}-\frac{x_2^2}{2\sigma_2^2}}$,
следователно
\[
  f_{Z,W}(z,w) = \frac1{2\pi\sigma_1\sigma_2}\,e^{-\frac{(z-w)^2}{2\sigma_1^2}-\frac{w^2}{2\sigma_2^2}}.
\]
Търсим маргиналната плътност на $Z=W_1+W_2$, интегрирайки по $w$.
Разгръщаме показателя (групирайки по степени на $w$):
\[
  -\frac{(z-w)^2}{2\sigma_1^2}-\frac{w^2}{2\sigma_2^2} = -\frac{w^2}{2\sigma_1^2}-\frac{w^2}{2\sigma_2^2}+\frac{wz}{\sigma_1^2}-\frac{z^2}{2\sigma_1^2}.
\]
Полагаме $\sigma^2:=\sigma_1^2+\sigma_2^2$, така че коефициентът пред
$w^2$ е $-\tfrac1{2\sigma_1^2}-\tfrac1{2\sigma_2^2}=-\tfrac{\sigma^2}
{2\sigma_1^2\sigma_2^2}$, и~допълваме до точен квадрат по общата
формула $-Aw^2+Bw+C=-A\big(w-\tfrac B{2A}\big)^2+\tfrac{B^2}{4A}+C$ с~$A=\tfrac{\sigma^2}{2\sigma_1^2\sigma_2^2}$, $B=\tfrac{z}{\sigma_1^2}$,
$C=-\tfrac{z^2}{2\sigma_1^2}$:
\[
  \frac B{2A} = \frac{z}{\sigma_1^2}\cdot\frac{\sigma_1^2\sigma_2^2}{\sigma^2} = \frac{\sigma_2^2z}{\sigma^2},
  \qquad
  \frac{B^2}{4A}+C = \frac{z^2\sigma_2^2}{2\sigma_1^2\sigma^2} - \frac{z^2}{2\sigma_1^2} = \frac{z^2}{2\sigma_1^2}\cdot\frac{\sigma_2^2-\sigma^2}{\sigma^2} = -\frac{z^2}{2\sigma^2}
\]
(последното — понеже $\sigma_2^2-\sigma^2=-\sigma_1^2$). Следователно
показателят се разпада точно на
\[
  -\frac{(z-w)^2}{2\sigma_1^2}-\frac{w^2}{2\sigma_2^2} = -\frac{\sigma^2}{2\sigma_1^2\sigma_2^2}\Big(w-\frac{\sigma_2^2}{\sigma^2}z\Big)^{\!2} \;-\; \frac{z^2}{2\sigma^2}.
\]
Извършваме смяната $v=\dfrac{\sigma}{\sigma_1\sigma_2}\Big(w-\dfrac{\sigma_2^2}{\sigma^2}z\Big)$
(т.е.\ $dw=\dfrac{\sigma_1\sigma_2}{\sigma}\,dv$), която превръща
експонентата в~$-\tfrac12v^2$:
\begin{align*}
  f_Z(z) &= \frac1{2\pi\sigma_1\sigma_2}\,e^{-\frac{z^2}{2\sigma^2}}\cdot\underbrace{\int_{-\infty}^{\infty} e^{-\frac12v^2}\,\frac{\sigma_1\sigma_2}{\sigma}\,dv}_{=\ \frac{\sigma_1\sigma_2}{\sigma}\sqrt{2\pi}\ \text{(нормировъчен интеграл на Гаус)}}\\
  &= \frac1{2\pi\sigma_1\sigma_2}\cdot\frac{\sigma_1\sigma_2}{\sigma}\sqrt{2\pi}\,e^{-\frac{z^2}{2\sigma^2}} = \frac1{\sqrt{2\pi}\,\sigma}\,e^{-\frac{z^2}{2\sigma^2}} \sim N(0,\sigma_1^2+\sigma_2^2). \qedbg
\end{align*}
\end{proofbox}

\section{Гама разпределение}

\begin{defbox}[~(Гама функция и~Гама разпределение)]
Нека $\alpha>0$. \textbf{(Ойлеровата) Гама функция} е
\[
  \Gamma(\alpha) := \int_0^{\infty} x^{\alpha-1}e^{-x}\,dx.
\]
Известни свойства: $\Gamma(1)=1$, $\Gamma(n{+}1)=n!$ за $n\in\mathbb
N$, и~рекурентната формула $\Gamma(\alpha{+}1)=\alpha\,\Gamma(\alpha)$
за $\alpha>0$. Нека $\alpha>0$, $\beta>0$. Казваме, че $X\sim
\Gamma(\alpha,\beta)$ (\textbf{Гама разпределение}), ако
\[
  f_X(x) = \begin{cases}\dfrac{\beta^\alpha}{\Gamma(\alpha)}\,x^{\alpha-1}e^{-\beta x}, & x>0\\ 0,&\text{иначе}\end{cases}.
\]
\end{defbox}

\begin{corbox}[~(частни случаи)]
При $\alpha=\beta=1$: $X\sim Exp(1)$. При $\alpha=1$, $\beta>0$
произволно: $f_X(x)=\dfrac{\beta}{\Gamma(1)}e^{-\beta x}=\beta
e^{-\beta x}$ за $x>0$, т.е.\ $\Gamma(1,\beta)=Exp(\beta)$.
\end{corbox}

\begin{thmbox}[]{Теорема}
Нека $X_i\sim\Gamma(\alpha_i,\beta)$, $i=1,\dots,n$ (\emph{еднакво}
$\beta$), независими в~съвкупност. Тогава
\[
  Y = \sum_{j=1}^n X_j \sim \Gamma(\alpha_1+\dots+\alpha_n,\ \beta).
\]
\end{thmbox}

\begin{proofbox}
Достатъчно е да докажем твърдението за $n=2$ (общият случай следва
по индукция). Нека $Y_1=X_1+X_2$, $Y_2=X_2$ (т.е.\ $X_1=Y_1-Y_2$,
$X_2=Y_2$), с~Якобиан $|J|=\det\begin{pmatrix}1&-1\\0&1\end{pmatrix}=1$.
Понеже $X_1\perp\!\!\!\perp X_2$:
\[
  f_X(x_1,x_2) = \frac{\beta^{\alpha_1+\alpha_2}}{\Gamma(\alpha_1)\Gamma(\alpha_2)}\,e^{-\beta(x_1+x_2)}x_1^{\alpha_1-1}x_2^{\alpha_2-1},
\]
и~след смяната на променливите:
\[
  f_{Y_1,Y_2}(y_1,y_2) = \frac{\beta^{\alpha_1+\alpha_2}}{\Gamma(\alpha_1)\Gamma(\alpha_2)}\,e^{-\beta y_1}(y_1-y_2)^{\alpha_1-1}y_2^{\alpha_2-1}, \qquad 0<y_2<y_1.
\]
Интегрираме по $y_2$, за да получим маргиналната плътност на
$Y_1=X_1+X_2$:
\[
  f_{Y_1}(y_1) = \frac{\beta^{\alpha_1+\alpha_2}}{\Gamma(\alpha_1)\Gamma(\alpha_2)}\,e^{-\beta y_1}\int_0^{y_1}(y_1-y_2)^{\alpha_1-1}y_2^{\alpha_2-1}\,dy_2.
\]
Полагаме $y_2=vy_1$ (т.е.\ $v\in(0,1)$), $dy_2=y_1\,dv$:
\[
  \int_0^{y_1}(y_1-y_2)^{\alpha_1-1}y_2^{\alpha_2-1}\,dy_2 = y_1^{\alpha_1+\alpha_2-1}\underbrace{\int_0^1 (1-v)^{\alpha_1-1}v^{\alpha_2-1}\,dv}_{=\ B(\alpha_1,\alpha_2)\ =\ \frac{\Gamma(\alpha_1)\Gamma(\alpha_2)}{\Gamma(\alpha_1+\alpha_2)}}
\]
(последният интеграл е класическата Ойлерова Бета-функция, за която
е известна връзката с~Гама функцията). Заместваме:
\[
  f_{Y_1}(y_1) = \frac{\beta^{\alpha_1+\alpha_2}}{\Gamma(\alpha_1)\Gamma(\alpha_2)}\,e^{-\beta y_1}\,y_1^{\alpha_1+\alpha_2-1}\cdot\frac{\Gamma(\alpha_1)\Gamma(\alpha_2)}{\Gamma(\alpha_1+\alpha_2)} = \frac{\beta^{\alpha_1+\alpha_2}}{\Gamma(\alpha_1+\alpha_2)}\,y_1^{\alpha_1+\alpha_2-1}e^{-\beta y_1},
\]
т.е.\ $Y_1\sim\Gamma(\alpha_1+\alpha_2,\beta)$. \qedbg
\end{proofbox}

\begin{corbox}[]
Ако $X_1,\dots,X_n$ са независими в~съвкупност, еднакво разпределени
$X_i\sim Exp(\lambda)$ (т.е.\ $X_i\sim\Gamma(1,\lambda)$), то
$\sum_{i=1}^n X_i \sim \Gamma(n,\lambda)$.
\end{corbox}

\begin{thmbox}[]{Теорема}
Ако $X\sim\Gamma(\alpha,\beta)$, то $\mathbb EX=\dfrac\alpha\beta$ и~$DX=\dfrac\alpha{\beta^2}$.
\end{thmbox}

\begin{proofbox}
Извършваме смяна на променливата $y=\beta x$ ($x=y/\beta$,
$dx=dy/\beta$), избрана така, че да разпознаем плътността на
$\Gamma(\alpha{+}1,\beta)$ (респ.\ $\Gamma(\alpha{+}2,\beta)$) под
интеграла:
\begin{align*}
  \mathbb EX &= \frac{\beta^\alpha}{\Gamma(\alpha)}\int_0^{\infty}x\cdot x^{\alpha-1}e^{-\beta x}\,dx = \frac{\beta^\alpha}{\Gamma(\alpha)}\int_0^{\infty}x^{\alpha}e^{-\beta x}\,dx.
\end{align*}
Разпознаваме под интеграла плътността на $\Gamma(\alpha{+}1,\beta)$
(с~изключение на нормиращата константа), т.е.\
$\int_0^{\infty}x^{\alpha}e^{-\beta x}\,dx = \dfrac{\Gamma(\alpha+1)}{\beta^{\alpha+1}}$
(понеже $\frac{\beta^{\alpha+1}}{\Gamma(\alpha+1)}x^{\alpha}e^{-\beta x}$ е плътност и~интегрира до 1). Следователно
\begin{align*}
  \mathbb EX = \frac{\beta^\alpha}{\Gamma(\alpha)}\cdot\frac{\Gamma(\alpha+1)}{\beta^{\alpha+1}} = \frac{\Gamma(\alpha+1)}{\beta\,\Gamma(\alpha)} = \frac{\alpha\,\Gamma(\alpha)}{\beta\,\Gamma(\alpha)} = \frac{\alpha}{\beta}.
\end{align*}
Аналогично, разпознавайки плътността на $\Gamma(\alpha{+}2,\beta)$:
\[
  \mathbb EX^2 = \frac{\beta^\alpha}{\Gamma(\alpha)}\int_0^{\infty}x^{\alpha+1}e^{-\beta x}\,dx = \frac{\Gamma(\alpha+2)}{\beta^2\,\Gamma(\alpha)} = \frac{(\alpha+1)\alpha\,\Gamma(\alpha)}{\beta^2\,\Gamma(\alpha)} = \frac{\alpha^2+\alpha}{\beta^2}
\]
(използвахме $\Gamma(\alpha{+}2)=(\alpha{+}1)\Gamma(\alpha{+}1)=
(\alpha{+}1)\alpha\,\Gamma(\alpha)$). Следователно
\[
  DX = \mathbb EX^2-(\mathbb EX)^2 = \frac{\alpha^2+\alpha}{\beta^2} - \frac{\alpha^2}{\beta^2} = \frac{\alpha}{\beta^2}. \qedbg
\]
\end{proofbox}

\section{\texorpdfstring{$\chi^2$}{Хи-квадрат} разпределение}

\begin{defbox}[~(\texorpdfstring{$\chi^2$}{хи-квадрат} разпределение)]
\textbf{Хи-квадрат разпределението} с~$n$ \emph{степени на свобода}
е частен случай на Гама разпределението: $X\sim\chi^2(n)$, ако
$X\sim\Gamma\big(\tfrac n2,\tfrac12\big)$, т.е.
\[
  f_X(x) = \begin{cases}\dfrac{(1/2)^{n/2}}{\Gamma(n/2)}\,x^{n/2-1}e^{-x/2}, & x>0\\ 0,&\text{иначе}\end{cases}.
\]
От формулите за Гама разпределение: $\mathbb EX = \dfrac{n/2}{1/2}=n$,
$DX = \dfrac{n/2}{(1/2)^2}=2n$.
\end{defbox}

\begin{thmbox}[]{Теорема}
Нека $X_1,\dots,X_n\sim N(0,1)$ са независими в~съвкупност. Тогава
\[
  Y = \sum_{i=1}^{n} X_i^2 \sim \chi^2(n).
\]
\end{thmbox}

\begin{proofbox}
Ако докажем, че $X_1^2\sim\chi^2(1)=\Gamma\big(\tfrac12,\tfrac12\big)$,
то по теоремата за сумата на независими Гама-разпределени величини
(с~еднакво $\beta=\tfrac12$): $\sum_{i=1}^n X_i^2 \sim
\Gamma\big(\tfrac n2,\tfrac12\big)=\chi^2(n)$.

Нека $H=X_1^2$ ($H$ не е монотонна функция на $X_1$ - параболата не
е инжективна - затова минаваме през функцията на разпределение
директно, вместо през общата теорема за смяна на променливите). За
$h>0$:
\[
  \mathbb P(H\le h) = \mathbb P(X_1^2\le h) = \mathbb P(-\sqrt h\le X_1\le\sqrt h) = 2\,\mathbb P(0\le X_1\le\sqrt h) = \frac2{\sqrt{2\pi}}\int_0^{\sqrt h} e^{-y^2/2}\,dy.
\]
Диференцираме директно по $h$ (правило на Лайбниц за диференциране
на интеграл с~граница, зависеща от параметъра), за да получим
плътността $f_H(h)=\tfrac{d}{dh}\mathbb P(H\le h)$:
\[
  f_H(h) = \frac2{\sqrt{2\pi}}\cdot e^{-h/2}\cdot\frac{d}{dh}\sqrt h = \frac2{\sqrt{2\pi}}\cdot e^{-h/2}\cdot\frac1{2\sqrt h} = \frac1{\sqrt{2\pi}}\,h^{-1/2}e^{-h/2}.
\]
Понеже $\Gamma(\tfrac12)=\sqrt\pi$, плътността на $\Gamma\big(\tfrac12,
\tfrac12\big)$ е точно $\dfrac{(1/2)^{1/2}}{\Gamma(1/2)}h^{-1/2}e^{-h/2}
= \dfrac1{\sqrt2\sqrt\pi}h^{-1/2}e^{-h/2} = \dfrac1{\sqrt{2\pi}}
h^{-1/2}e^{-h/2}$, което съвпада с~$f_H(h)$. Следователно
$H=X_1^2\sim\Gamma\big(\tfrac12,\tfrac12\big)=\chi^2(1)$. \qedbg
\end{proofbox}

\section{\texorpdfstring{$t$}{t}-разпределение (на Стюдънт)}

\begin{defbox}[~(\texorpdfstring{$t$}{t}-разпределение)]
Нека $n\ge1$. Казваме, че $X\sim t(n)$ с~$n$ степени на свобода, ако
\[
  X = \frac{Z}{\sqrt{Y/n}}, \qquad Z\perp\!\!\!\perp Y,\quad Z\sim N(0,1),\quad Y\sim\chi^2(n).
\]
\end{defbox}
